BVRIJHKLMN photometry of 2 stars from the constellation Cassiopeia is given in
Table 5. For both stars it is believed that their light is affected by
extinction by diffuse Interstellar Medium (ISM) only. Assuming that the
observation is done from outside the atmosphere.

\begin{description}
\item[a)] Using the data given in Tables 5 to 9, plot $E_{X-V}/E_{B-V}$ as a
  function of $1/\lambda_X$ for filters B, V, R, I, J, H, K, L, M, N for both
  stars. Fit approximate curves by eye (in particular, note that
  $E_{X-V}/E_{B-V} \sim \mathit{const.}$ as $1/\lambda_X \to 0$). X is each
  band in the photometric system.

  $E_{B-V}$ is the colour excess.
\item[b)] Using the graphs obtained in a), estimate $R_V$ and $R_R$ for each
  star.
  \[ R_V = \frac{A_V}{E_{B-V}} \quad \text{and} \quad
     R_R = \frac{A_R}{E_{R-I}} \]
  ($A_V$ is the absorption in V).
\end{description}

Now apply these results in order to derive a distance estimate for IC 342, a
spiral galaxy in Cassiopeia obscured by Milky Way. You should assume that the
properties of the ISM in IC 342 are similar to those of the ISM in our Galaxy.

\begin{description}
\item[c)] Using the period-magnitude diagrams for 20 Cepheids from IC 342
  (Figures 2 and 3) and assuming the period-luminosity relations:
  \[ \langle M_R \rangle = -2.91 \left( \log\left(\frac{P}{\mathrm{day}}\right)
     - 1 \right) - 4.04 \quad \text{and} \quad
     \langle M_I \rangle = -3.00 \left( \log\left(\frac{P}{\mathrm{day}}\right)
     - 1 \right) - 4.06 \]
  where $\langle M_R \rangle$ and $\langle M_I \rangle$ are the mean absolute
  magnitudes in filters R and I, find $A_R$ for objects in IC 342. Find the
  distance to IC 342.
\end{description}

\medskip
\noindent Table 5: BVRIJHKLMN photometry of two stars in Cassiopeia

\begin{center}
\begin{tabular}{lccccccccccc}
\hline
Star & MK & $B$ & $V$ & $R$ & $I$ & $J$ & $H$ & $K$ & $L$ & $M$ & $N$ \\
 & class & mag & mag & mag & mag & mag & mag & mag & mag & mag & mag \\
\hline
HD 4817 & K3Iab & 8.08 & 6.18 & 4.73 & 3.64 & 2.76 & 1.86 & 1.54 & 1.32 &
  1.59 & -- \\
HD 11092 & K4II & 8.66 & 6.57 & -- & -- & 3.10 & 2.14 & 1.63 & 1.41 & 1.65 &
  1.44 \\
\hline
\end{tabular}
\end{center}

\medskip
\noindent Table 6: $(B-V)_0$ intrinsic colours for selected sp.\ types and
luminosity classes

\begin{center}
\begin{tabular}{ccc}
\hline
 & \multicolumn{2}{c}{$(B-V)_0$ [mag]} \\
 & II & Iab / Ia \\
\hline
F0 & -- & 0.15 \\
G0 & 0.73 & 0.82 \\
K0 & 1.06 & 1.18 \\
K3 & 1.40 & 1.42 \\
K4 & 1.42 & 1.50 \\
\hline
\end{tabular}
\end{center}

\medskip
\noindent Table 7: Infrared intrinsic colours for selected sp.\ types of
supergiant stars

\begin{center}
\begin{tabular}{ccccccccc}
\hline
 & $(V-R)_0$ & $(V-I)_0$ & $(V-J)_0$ & $(V-H)_0$ & $(V-K)_0$ & $(V-L)_0$ &
   $(V-M)_0$ & $(V-N)_0$ \\
 & mag & mag & mag & mag & mag & mag & mag & mag \\
\hline
F0 & 0.20 & 0.31 & 0.36 & 0.51 & 0.60 & 0.64 & 0.65 & 0.82 \\
G0 & 0.55 & 0.90 & 1.14 & 1.52 & 1.71 & 1.72 & 1.72 & 1.98 \\
K0 & 0.95 & 1.59 & 2.01 & 2.64 & 2.80 & 2.87 & 2.79 & 3.14 \\
K3 & 1.13 & 1.96 & 2.41 & 3.14 & 3.25 & 3.39 & 3.25 & 3.63 \\
K4 & 1.20 & 2.13 & 2.59 & 3.37 & 3.44 & 3.62 & 3.46 & 3.84 \\
\hline
\end{tabular}
\end{center}

\medskip
\noindent Table 8: Infrared intrinsic colours for selected sp.\ types of
giant stars

\begin{center}
\begin{tabular}{ccccccccc}
\hline
 & $(V-R)_0$ & $(V-I)_0$ & $(V-J)_0$ & $(V-H)_0$ & $(V-K)_0$ & $(V-L)_0$ &
   $(V-M)_0$ & $(V-N)_0$ \\
 & mag & mag & mag & mag & mag & mag & mag & mag \\
\hline
K0 & 0.60 & 1.03 & 1.23 & 1.72 & 1.94 & 1.97 & 1.90 & 1.92 \\
K3 & 0.86 & 1.39 & 1.84 & 2.40 & 2.69 & 2.82 & 2.70 & 2.73 \\
K4 & 0.96 & 1.61 & 2.16 & 2.77 & 3.05 & 3.22 & 3.08 & 3.02 \\
\hline
\end{tabular}
\end{center}

\medskip
\noindent Table 9: Effective wavelengths of selected photometric filters

\begin{center}
\begin{tabular}{ccccccccccc}
\hline
Filter & B & V & R & I & J & H & K & L & M & N \\
$\lambda_{\mathrm{F}}$/nm & 450 & 555 & 670 & 870 & 1200 & 1620 & 2200 &
  3500 & 5000 & 9000 \\
\hline
\end{tabular}
\end{center}

% Fig. 2: <R> is the mean apparent magnitude in filter R
\ioaafig{2015_DA_D2-f5.pdf}
% Fig. 3: <I> is the mean apparent magnitude in filter I
\ioaafig{2015_DA_D2-f6.pdf}
